\documentclass{ximera}

\begin{document}
	\author{Stitz-Zeager}
	\xmtitle{Equations and Inequalities involving Logarithmic Functions}{}


\mfpicnumber{1}

\opengraphsfile{LogarithmicEquationsandInequalities}

\setcounter{footnote}{0}

\label{LogarithmicEquationsandInequalities}

In Section \ref{ExponentialEquationsandInequalities} we solved equations and inequalities involving exponential functions using one of two basic strategies.  We now turn our attention to equations and inequalities involving logarithmic functions, and not surprisingly, there are two basic strategies to choose from.  



For example, per Theorem \ref{explogsonetoone}, the \textit{only} solution to $\log_{2}(x) = \log_{2}(5)$ is $x=5$.   Now consider $\log_{2}(x) = 3$.  To use Theorem \ref{explogsonetoone}, we need to rewrite $3$ as a logarithm base $2$.   Theorem \ref{invpropslogs} gives us $3 = \log_{2}\left(2^{3}\right) = \log_{2}(8)$.  Hence, $\log_{2}(x) = 3$ is equivalent to  $\log_{2}(x) =  \log_{2}(8)$ so that $x = 8$. 

A second approach to solving  $\log_{2}(x) = 3$ us to apply the corresponding exponential function, $f(x) = 2^x$ to both sides:  $2^{\log_{2}(x)} = 2^{3}$ so $x = 2^3 = 8$.

A third approach to solving  $\log_{2}(x) = 3$ is to use  Theorem \ref{invpropslogs} to rewrite $\log_{2}(x) = 3$ as $2^{3} = x$, so  $x=8$.   

In the grand scheme of things, all three approaches we have presented to solve $\log_{2}(x) = 3$ are mathematically equivalent, so we opt to choose the last approach in our summary below.

%% \colorbox{ResultColor}{\bbm

\centerline{\textbf{Steps for Solving an Equation involving Logarithmic Functions}} \index{logarithm ! solving equations with} 

\begin{enumerate}

\item  Isolate the logarithmic function.

\item  \begin{enumerate}

\item  If convenient, express both sides as logs with the same base and equate arguments.

\item  Otherwise, rewrite the log equation as an exponential equation.


\end{enumerate}

\end{enumerate}



%% \ebm}



\begin{example}  \label{LogEqnsEx1} Solve the following equations.  Check your solutions graphically using a graphing utility.

\begin{enumerate}

\item  $\log_{117}(1-3x) = \log_{117}\left(x^2-3\right)$

\item  $2 - \ln(t-3) = 1$

\item  $\log_{6}(x+4) + \log_{6}(3-x) = 1$

\item  $\log_{7}(1-2t) = 1 - \log_{7}(3-t)$

\item  $\log_{2}(x+3) = \log_{2}(6-x)+3$

\item  $1 + 2 \log_{4}(t+1) = 2 \log_{2}(t)$

\end{enumerate}

\begin{explanation}

\begin{enumerate}

\item  Since we have the same base on both sides of the equation $\log_{117}(1-3x) = \log_{117}\left(x^2-3\right)$, we equate the arguments (what's inside) of the logs to get $1-3x = x^2-3$.  Solving $x^2+3x-4 = 0$ gives $x=-4$ and $x=1$.  

Using desmos, we graph both $y = (x)$ and $y = g(x)$ and find that these graphs intersect only at $x=-4$.


% https://www.desmos.com/calculator/xlwl2vsugi
\begin{center}
    \desmos{xlwl2vsugi}{800}{600}
\end{center}

To see what happened to the solution $x=1$, we substitute it into our original equation to obtain  $\log_{117}(-2) =  \log_{117}(-2)$.  While these expressions look identical, neither is a real number,\footnote{They do, however, represent the same \textbf{family} of complex numbers.  We refer the reader to a course in Complex Variables.} which means $x=1$ is not in the domain of the original equation, and is not a solution.  


\item  To solve  $2 - \ln(t-3) = 1$, we first isolate the logarithm and get $\ln(t-3) = 1$. Rewriting $\ln(t-3) = 1$ as an exponential equation, we get is $e^{1} = t-3$, so $t =e+3$. 

Desmos shows the graphs of $f(t) = 2 - \ln(t-3)$ and $g(t) = 1$ intersect at $t \approx   5.718 \approx e+3$.

% https://www.desmos.com/calculator/la8yx4pevj
\begin{center}
\desmos{la8yx4pevj}{800}{600}
\end{center}


\item We  start solving $\log_{6}(x+4) + \log_{6}(3-x) = 1$ by using the Product Rule for logarithms to rewrite the equation as  $\log_{6}\left[(x+4)(3-x)\right] = 1$.  

Rewriting as an exponential equation gives $6^{1} = (x+4)(3-x)$ which reduces to $x^2+x-6 = 0$.  We get two solutions: $x=-3$ and $x=2$.   

We graph $y=f(x)$ and $y=g(x)$ using desmos and see the graphs intersect twice, at $x=-3$ and $x=2$, as required.

% https://www.desmos.com/calculator/u2zjqnhnnz
\begin{center}
    \desmos{u2zjqnhnnz}{800}{600}
\end{center}

\item  Taking a cue from the previous problem, we begin solving $\log_{7}(1-2t) = 1 - \log_{7}(3-t)$ by first collecting the logarithms on the same side, $\log_{7}(1-2t) +  \log_{7}(3-t) = 1$, and then using the Product Rule to get $\log_{7}[(1-2t)(3-t)] = 1$.  

Rewriting  as an exponential equation gives $7^{1} = (1-2t)(3-t)$ which gives the quadratic equation $2t^2-7t-4=0$.  Solving, we find  $t = -\frac{1}{2}$ and $t=4$.  

Once again, we use the change of base formula and find the graphs of  $y = f(t) = \frac{\ln(1-2t)}{\ln(7)}$ and $y=g(t) = 1 - \frac{\ln(3-t)}{\ln(7)}$ intersect only at $t=-\frac{1}{2}$.  

Checking $t=4$ in the original equation produces $\log_{7}(-7) = 1 - \log_{7}(-1)$, showing $t=4$ is not in the domain of $f$ nor $g$.

% https://www.desmos.com/calculator/vuoatn2epd
\begin{center}
\desmos{vuoatn2epd}{800}{600}
\end{center}

\item Our first step in solving  $\log_{2}(x+3) = \log_{2}(6-x)+3$ is to gather the logarithms to one side of of the equation: $\log_{2}(x+3) - \log_{2}(6-x) = 3$.

The  Quotient Rule gives $\log_{2}\left(\frac{x+3}{6-x}\right) = 3$ which, as an exponential equation is $2^{3} = \frac{x+3}{6-x}$. 

Clearing denominators, we get $8(6-x) = x+3$, which reduces to  $x = 5$. 

Graphing $f$ and $g$ using desmos, we find the graphs intersect at $x=5$.

% https://www.desmos.com/calculator/blhn7wqp22
\begin{center}
    \desmos{blhn7wqp22}{800}{600}
\end{center}


\item Our first step in solving $1 + 2 \log_{4}(t+1) = 2 \log_{2}(t)$ is to gather the logs on one side of the equation.  We obtain  $1 = 2 \log_{2}(t) - 2 \log_{4}(t+1)$ but find we need a common base to combine the logs.

Since $4$ is a power of $2$, we use change of base to convert  $\log_{4}(t+1) = \frac{\log_{2}(t+1)}{\log_{2}(4)} = \frac{1}{2} \log_{2}(t+1)$.  Hence, our original equation becomes  

\[ \begin{array}{rclr}

1 & = & 2 \log_{2}(t) - 2 \left(\frac{1}{2} \log_{2}(t+1)\right) & \\  
1 &= & 2\log_{2}(t) - \log_{2}(t+1) & \\  
1 & = & \log_{2}\left(t^2\right) - \log_{2}(t+1) & \text{Power Rule} \\  
1 & = & \log_{2}\left( \dfrac{t^{2}}{t+1}\right) & \text{Quotient Rule} \\ \end{array}\]

Rewriting $1 = \log_{2}\left( \frac{t^{2}}{t+1}\right)$ in exponential form gives  $ \frac{t^{2}}{t+1} = 2$ or $t^2 -2t-2 = 0$.  Using the quadratic formula, we obtain  $t = 1 \pm \sqrt{3}$.

We use desmos to graph $y = f(t)$ and $y = g(t)$ below. We see the graphs intersect only at $t \approx 2.732 \approx 1 + \sqrt{3}$.  


% https://www.desmos.com/calculator/befiv6qmbt
\begin{center}
\desmos{befiv6qmbt}{800}{600}
\end{center}

Note the solution $t = 1 - \sqrt{3} < 0$ Hence if substituted into the original equation, the term $2 \log_{2}\left(1 - \sqrt{3}\right)$ is undefined, which explains why the graphs intersect only once.


\end{enumerate}
\end{explanation}
\end{example}

If nothing else,  Example \ref{LogEqnsEx1} demonstrates the importance of checking for extraneous solutions\footnote{Recall that an extraneous solution is an answer obtained analytically which does not satisfy the original equation.} when solving equations involving logarithms.  Even though we checked our answers graphically, extraneous solutions are easy to spot:  any supposed solution which causes the argument of a logarithm to be negative must be discarded.  

While identifying extraneous solutions is important, it is equally important to understand which machinations create the opportunity for extraneous solutions to appear.  In the case of Example \ref{LogEqnsEx1}, extraneous solutions, by and large, result from using the Power, Product, or Quotient Rules.  We encourage the reader to take the time to track each extraneous solution found in  Example \ref{LogEqnsEx1}  backwards through the solution process to see at precisely which step it fails to be a solution.

As with the equations in Example \ref{expeqnsex1}, much can be learned from checking all of the answers in Example \ref{LogEqnsEx1} analytically.  We leave this to the reader and turn our attention to inequalities involving logarithmic functions.  Since logarithmic functions are continuous on their domains, we can use sign diagrams.  

\begin{example}  Solve the following inequalities.  Check your answer graphically using a calculator.
\label{logineq}

\begin{enumerate}

\item  $\dfrac{1}{\ln(x)+1} \leq 1$

\item  $\left(\log_{2}(x)\right)^2 < 2 \log_{2}(x) + 3$

\item  $t \log(t+1) \geq t$

\end{enumerate}


\begin{explanation}

\begin{enumerate}

\item  We start solving $\frac{1}{\ln(x)+1} \leq 1$ by getting $0$ on one side of the inequality: $\frac{1}{\ln(x)+1}  - 1 \leq 0$.  

Getting a common denominator yields $\frac{1}{\ln(x)+1}  - \frac{\ln(x)+1}{\ln(x)+1} \leq 0$ which reduces to $\frac{-\ln(x)}{\ln(x)+1} \leq 0$, or $ \frac{\ln(x)}{\ln(x)+1} \geq 0$.  

We define $r(x) = \frac{\ln(x)}{\ln(x)+1}$ and set about finding the domain and the zeros of $r$.  Due to the appearance of the term $\ln(x)$, we require  $x > 0$.  In order to keep the denominator away from zero, we solve $\ln(x)+1 = 0$ so $\ln(x) = -1$, so $x = e^{-1} = \frac{1}{e}$.  Hence, the domain of $r$ is $\left(0, \frac{1}{e}\right) \cup \left(\frac{1}{e}, \infty\right)$. 

To find the zeros of $r$, we set $r(x) = \frac{\ln(x)}{\ln(x)+1} = 0$ so that $\ln(x) = 0$, and we find $x = e^{0} = 1$.  


In order to determine test values for $r$ without resorting to the calculator, we need to find numbers between $0$, $\frac{1}{e}$, and $1$ which have a base of $e$.  Since $e \approx 2.718 > 1$, $0 < \frac{1}{e^2} < \frac{1}{e} < \frac{1}{\sqrt{e}} < 1 < e$.  



To determine the sign of $r\left( \frac{1}{e^2} \right)$, note $\ln\left(\frac{1}{e^2}\right) = \ln\left(e^{-2}\right) = -2$. Hence, $r\left( \frac{1}{e^2} \right) = \frac{-2}{-2+1} = 2 > 0$.  The rest of the test values are determined similarly.   


\begin{image}
% \input{PowerFunctions_pic26.tex}
\begin{tikzpicture}[x=20pt,y=20pt]
  \draw[-{Latex[length=5pt,width=3pt]}, line width=1.25pt] (-4,0) -- (2,0);
  \draw (-2,-0.12) -- (-2,0.12);
  \draw (0,-0.12) -- (0,0.12);
  \node at (-3, 0.5) {$(+)$};
  \node at (-2,-0.5) {$1/e$};
  \node at (-2,0.5) {\textinterrobang};
  \node at (-1,0.5) {$(-)$};
  \node at (0,-0.5) {$1$};
  \node at (0,0.5) {$0$};
  \node at (-4,-0.5) {$0$};
  \node at (1,0.5) {$(+)$};
  \draw[fill=white, draw=black] (-4,0) circle (0.1cm);
\end{tikzpicture}
\end{image}

From our sign diagram, we find $r(x) \geq 0$ on $\left(0, \frac{1}{e}\right) \cup [1, \infty)$, which is our solution. 


Graphing $f(x) =  \frac{1}{\ln(x)+1}$ and $g(x) = 1$ using desmos, we see the graph of $f$ is below the graph of $g$ on these intervals, and that the graphs intersect at $x=1$.

 % https://www.desmos.com/calculator/ytoau4uepr
\begin{center}
    \desmos{ytoau4uepr}{800}{600}
\end{center}

\item  Moving all of the nonzero terms of  $\left(\log_{2}(x)\right)^2 < 2 \log_{2}(x) + 3$ to one side of the inequality in order to make use of a sign diagram, we have $\left(\log_{2}(x)\right)^2 - 2 \log_{2}(x) - 3 < 0$. 

Defining $r(x) = \left(\log_{2}(x)\right)^2 - 2 \log_{2}(x) - 3$, we get the domain of $r$ is $(0, \infty)$, due to the presence of the logarithm.  To find the zeros of $r$, we set $r(x) =\left(\log_{2}(x)\right)^2 - 2 \log_{2}(x) - 3= 0$ which we identify as a  `quadratic in disguise.'  

Setting $u = \log_{2}(x)$, our equation becomes $u^2-2u-3 = 0$.  Factoring   gives us $u=-1$ and $u=3$.  Since $u = \log_{2}(x)$, we get $\log_{2}(x) = -1$, or $x = 2^{-1} = \frac{1}{2}$, and $\log_{2}(x) = 3$, which gives  $x = 2^{3} = 8$.  


We use test values which are powers of $2$: $0 < \frac{1}{4} < \frac{1}{2} < 1 < 8 < 16$ to create the sign diagram below. 


\begin{image}
% \input{PowerFunctions_pic26.tex}
\begin{tikzpicture}[x=20pt,y=20pt]
  \draw[-{Latex[length=5pt,width=3pt]}, line width=1.25pt] (-4,0) -- (2,0);
  \draw (-2,-0.12) -- (-2,0.12);
  \draw (0,-0.12) -- (0,0.12);
  \node at (-3, 0.5) {$(+)$};
  \node at (-2,-0.5) {$1/2$};
  \node at (-2,0.5) {$0$};
  \node at (-1,0.5) {$(-)$};
  \node at (0,-0.5) {$8$};
    \node at (0,0.5) {$0$};
  \node at (-4,-0.5) {$0$};
  \node at (1,0.5) {$(+)$};
  \draw[fill=white, draw=black] (-4,0) circle (0.1cm);
\end{tikzpicture}
\end{image}


From our sign diagram, we see $r(x)< 0$, which corresponds to our solution,  on $\left(\frac{1}{2}, 8 \right)$. 


 < 

Using desmos, we graph  $f(x)= \left(\log_{2}(x)\right)^2$ and $g(x) = 2 \log_{2}(x) + 3$.  We find the graph of $f$ is below the graph of $g$ on $\left(\frac{1}{2}, 8 \right)$.



% https://www.desmos.com/calculator/vdgpnywilj
\begin{center}
    \desmos{vdgpnywilj}{800}{600}
\end{center}


\item  We begin to solve $t \log(t+1) \geq t$ by subtracting $t$ from both sides to get $t \log(t+1)  - t \geq 0$.  

We define $r(t) = t \log(t+1)  - t $ and due to the presence of the logarithm, we require $t > -1$.  

To find the zeros of $r$, we set $r(t) = t \log(t+1)  - t = 0$.  Factoring, we get $t \left(\log(t+1) - 1\right) = 0$, which gives $t=0$ or $\log(t+1) - 1=0$.  

From $\log(t+1) - 1=0$ we get   $\log(t+1) = 1$, which we  rewrite as  $t+1 = 10^{1}$. Hence,   $t = 9$.  


We select test values $t$ so that $t+1$ is a power of $10$. Using $-1 < -0.9 < 0 < \sqrt{10} -1 < 9 < 99$, we get the sign diagram below,

\begin{image}
% \input{PowerFunctions_pic26.tex}
\begin{tikzpicture}[x=20pt,y=20pt]
  \draw[-{Latex[length=5pt,width=3pt]}, line width=1.25pt] (-4,0) -- (2,0);
  \draw (-2,-0.12) -- (-2,0.12);
  \draw (0,-0.12) -- (0,0.12);
  \node at (-3, 0.5) {$(+)$};
  \node at (-2,-0.5) {$0$};
  \node at (-2,0.5) {$0$};
  \node at (-1,0.5) {$(-)$};
  \node at (0,-0.5) {$9$};
  \node at (0,0.5) {$0$};
  \node at (-4,-0.5) {$-1$};
  \node at (1,0.5) {$(+)$};
  \draw[fill=white, draw=black] (-4,0) circle (0.1cm);
\end{tikzpicture}
\end{image}

Our sign diagram gives the solution as $(-1,0] \cup [9, \infty)$, which we check graphically using desmos.  

We find the graphs of  $y= f(t) = t \log(t+1)$  and $y=g(t) = t$ intersect at $t=0$ and $t=9$ with the graph of $f$ above the graph of $g$ on the given solution intervals.

% https://www.desmos.com/calculator/nsc8divdxu
\begin{center}
    \desmos{nsc8divdxu}{800}{600}
\end{center}

\end{enumerate}
\end{explanation}
\end{example}



Our next example revisits the concept of pH first seen in Exercise \ref{pHexercise} in Section \ref{LogarithmicFunctions}.  


\begin{example}

In order to successfully breed Ippizuti fish the pH of a freshwater tank must be at least 7.8 but can be no more than 8.5.  Determine the corresponding range of hydrogen ion concentration, and check your answer using a calculator.



\begin{explanation}  Recall from Exercise \ref{pHexercise} in Section \ref{LogarithmicFunctions} that $\text{pH} = -\log[\text{H}^{+}]$ where $[\text{H}^{+}]$ is the hydrogen ion concentration in moles per liter.  

We require $7.8 \leq -\log[\text{H}^{+}] \leq 8.5$ or $-8.5 \leq   \log[\text{H}^{+}] \leq -7.8$.  One way to proceed is to break this compound inequality into two inequalities, solve each using a sign diagram, and take the intersection of the solution sets.\footnote{Refer to page \pageref{intersectionunion} for a discussion of what this means.}  

On the other hand, we take advantage of the fact that $F(x) = 10^{x}$ is an increasing function, meaning that if $a \leq b \leq c$, then $10^{a} \leq 10^{b} \leq 10^{c}$.  This property allows us to solve our inequality in one step: from $-8.5 \leq   \log[\text{H}^{+}] \leq -7.8$, we get   $10^{-8.5} \leq   10^{\log[\text{H}^{+}]} \leq 10^{-7.8}$, so our solution is $ \leq   [\text{H}^{+}] \leq 10^{-7.8}$.  (Your Chemistry professor may want the answer written as $3.16 \times 10^{-9} \leq [\text{H}^{+}] \leq 1.58 \times 10^{-8}$.)  Using interval notation, our answer is $\left[10^{-8.5}, 10^{-7.8}\right]$.   

After \textit{very} carefully adjusting the viewing window on desmos, we see the graph of $f(x) = -\log(x)$ lies between the lines $y = 7.8$ and $y = 8.5$ on the interval $[3.162 \times 10^{-9}, 1.5849 \times 10^{-8}]$.

% https://www.desmos.com/calculator/dhebz5kshi
\begin{center}
    \desmos{dhebz5kshi}{800}{600}
\end{center}
\end{explanation}
\end{example}



We close this section by finding an inverse of a one-to-one function which involves logarithms.

\begin{example}  \label{logfracinverse} The function $f(x) = \dfrac{\log(x)}{1-\log(x)}$ is one-to-one. 

\begin{enumerate}

\item  Find a formula for $f^{-1}(x)$ and check your answer graphically using a graphing utility.

\item Solve  $\dfrac{\log(x)}{1-\log(x)} = 1$

\end{enumerate}

\begin{explanation} \begin{enumerate} \item  We first write $y=f(x)$ then interchange the $x$ and $y$ and solve for $y$.

\[ \begin{array}{rclr}
y & = & f(x) & \\ 
y  & = & \dfrac{\log(x)}{1-\log(x)} & \\ 
x  & = & \dfrac{\log(y)}{1-\log(y)} & \text{Interchange $x$ and $y$.}\\ 
x\left(1-\log(y)\right) & = & \log(y) & \\ 
x - x\log(y)  & = & \log(y) & \\ 
x & = & x \log(y) + \log(y) & \\ 
x & = & (x+1) \log(y) & \\ 
\dfrac{x}{x+1}  & = & \log(y) & \\ 
y & = & 10^{\frac{x}{x+1}} & \text{Rewrite as an exponential equation.}\\

\end{array}\]



We have $f^{-1}(x) = 10^{\frac{x}{x+1}}$.  Graphing $f$ and $f^{-1}$ using desmos shows the required symmetry about the line $y=x$.

% https://www.desmos.com/calculator/gett0ddvxe
\begin{center}
\desmos{gett0ddvxe}{800}{600}
\end{center}

\item Recognizing   $\frac{\log(x)}{1-\log(x)} = 1$ as $f(x) = 1$, we have $x = f^{-1}(1) = 10^{\frac{1}{1+1}} = 10^{\frac{1}{2}} = \sqrt{10}$.

To check our answer algebraically,  first recall  $\log(\sqrt{10}) = \log_{10}(\sqrt{10})$.  Next, we know $\sqrt{10} = 10^{\frac{1}{2}}$.  Hence, $\log_{10} \left(10^{\frac{1}{2}} \right) = \frac{1}{2} = 0.5$.  It follows that $\frac{\log(\sqrt{10})}{1-\log(\sqrt{10})} = \frac{0.5}{1-0.5} = \frac{0.5}{0.5} = 1$, as required.  

\end{enumerate}
\end{explanation}

\end{example}


Our last example uses the tools from this section along with Section \ref{AppDerivatives}.

\begin{example}\label{logcurvesketchex} Let  $f(x) = x^2 \ln(x)$.

\begin{enumerate}

\item  Find the domain of $f$.

\item  Find the $x$-intercepts, if any.

\item Explain why $\lim_{x \rightarrow 0^{+}} x^2 \ln(x)$ results in an indeterminate form.

Use a table of values  to approximate $\lim_{x \rightarrow 0^{+}} x^2 \ln(x)$.  What does your answer mean graphically? 

\item  Find  $\lim_{x \rightarrow \infty} x^2 \ln(x)$ 

\item  Given $f'(x) = 2x \ln(x) + x$, find the open intervals over which $f$ is increasing and decreasing.

\item  Find the local extrema.

\item  Given $f''(x) = 2 \ln(x) + 3$, find the open intervals over which the graph of $f$ is concave up or concave down.

\item Find the inflection points.

\item  Check your answers using a graphing utility.

\end{enumerate} 

\begin{explanation}

\begin{enumerate}

\item  Owing to the presence of the `$\ln(x)$' factor, we have the domain of $f$ is $(0, \infty)$.

\item  To find the $x$-intercepts, we set $f(x) = x^2 \ln(x) = 0$.  This gives $x^2 = 0$ or $\ln(x) = 0$, so $x = 0$ or $x = 1$.  Since $x=0$ is not in the domain of $f$, our only solution is $x = 1$. We have just one $x$-intercept, $(1,0)$.

\item   To analyze $\lim_{x \rightarrow 0^{+}} x^2 \ln(x)$, we note that as $x \rightarrow 0^{+}$, $x^2 \rightarrow 0$ but $\ln(x) \rightarrow -\infty$.  Hence we obtain the indeterminate form\footnote{See the remarks following Example \ref{exponentialcurvesketchingex}.} `$0 \cdot (- \infty)$.'  Using desmos, we make a table of values of $f(x)$ as $x \rightarrow 0^{+}$.  Both the table and the graph suggest that  $\lim_{x \rightarrow 0^{+}} x^2 \ln(x) = 0$.  Since $x = 0$ is not in the domain of $f$, there is a hole in the graph at $(0,0)$.

% https://www.desmos.com/calculator/nctktbdghr
\begin{center}
\desmos{nctktbdghr}{800}{600}
\end{center}

\item  To determine the intervals over which $f$ is increasing and decreasing, we make a sign diagram for $f'(x) = 2x \ln(x) + x$.  Solving $f'(x) = 2x \ln(x) + x = 0$ we factor: $x(2\ln(x) + 1) = 0$.  Since the domain of $f$ is $(0,
\infty)$, we focus on the factor $2\ln(x) + 1 = 0$.  We get $\ln(x) = -\frac{1}{2}$ so $x = e^{-\frac{1}{2}}$.  

When choosing test values, we could opt to find a decimal approximation for $e^{-\frac{1}{2}}$ or we could use more `log-friendly' values.  In this case we note that $e^{-1} < e^{-\frac{1}{2}} < e^{0} = 1$ so we choose $e^{-1}$ and $1$ as our test values.

We find $f'\left(e^{-1}\right) = 2 \ln\left(e^{-1}\right) + 1 = 2(-1) + 1 = -1 <0$ and $f'(1) = 2 \ln(1) + 1 = 1 > 0$.  Our sign diagram for $f'(x)$ is below.

\begin{image}
\begin{tikzpicture}[>=Latex, every node/.style={font=\normalsize}]

\def\crit{1}
\def\testL{0}
\def\testR{2}
\def\xmin{-1}
\def\xmax{3.2}

\draw[->] (\xmin,0) -- (\xmax,0);
\draw (-1,0.08) -- (-1,-0.08);
\draw (\crit,0.08) -- (\crit,-0.08);
\node[below] at (\crit,-0.12) {$e^{-1/2}$};
\node[below] at (-1,-0.12) {$0$};
\node[above] at (\crit,0.15) {$0$};

\node[above] at ({(\xmin+\crit)/2},0.15) {$(-)$};
\node[above] at ({(\crit+\xmax-0.6)/2},0.15) {$(+)$};
\node[right] at (\xmax+0.15,0.55) {$f'(x)$};

\draw[->, gray] (\testL,-0.9) -- (\testL,-0.15);
\draw[->, gray] (\testR,-0.9) -- (\testR,-0.15);
\node[right] at (\xmax+0.15,-0.15) {$x$};
\node[below] at (\testL,-0.9) {$e^{-1}$};
\node[below] at (\testR,-0.9) {$1$};

\end{tikzpicture}

\end{image}

We interpret the sign diagram using Theorem \ref{firstderivatveandgraphs} as follows:

\begin{image}
\begin{tikzpicture}[>=Latex, every node/.style={font=\normalsize}]

\def\crit{1}
\def\testL{0}
\def\testR{2}
\def\xmin{-1}
\def\xmax{3.2}

\draw[->] (\xmin,0) -- (\xmax,0);
\draw (-1,0.08) -- (-1,-0.08);
\draw (\crit,0.08) -- (\crit,-0.08);
\node[below] at (\crit,-0.12) {$e^{-1/2}$};
\node[below] at (-1,-0.12) {$0$};
\node[above] at (\crit,0.15) {$\rightarrow$};

\node[above] at ({(\xmin+\crit)/2},0.15) {$\searrow$};
\node[above] at ({(\crit+\xmax-0.6)/2},0.15) {$\nearrow$};
\node[right] at (\xmax+0.15,0.55) {$f(x)$};

% \draw[->, gray] (\testL,-0.9) -- (\testL,-0.15);
% \draw[->, gray] (\testR,-0.9) -- (\testR,-0.15);
\node[right] at (\xmax+0.15,-0.15) {$x$};
% \node[below] at (\testL,-0.9) {$e^{-1}$};


\end{tikzpicture}

\end{image}


We find $f$ is decreasing on $\left(0, e^{-\frac{1}{2}} \right)$ and increasing on $\left( e^{-\frac{1}{2}}, \infty\right)$.

\item  We see $f$ has a local (and in this case, global) minimum when $x = e^{-\frac{1}{2}}$.  The minimum value is $f\left(e^{-\frac{1}{2}}\right) = \left(e^{-\frac{1}{2}}\right)^2 \ln\left(e^{-\frac{1}{2}}\right) = e^{-1} \left(-\frac{1}{2}\right) = -\frac{e^{-1}}{2}$. The local minimum is $\left( e^{-\frac{1}{2}}, -\frac{e^{-1}}{2} \right)$.



\item  To find the intervals over which the graph of $f$ is concave up and concave down, we make a sign diagram for $f''(x) = 2\ln(x) + 3$.  Solving $f''(x) = 2 \ln(x) + 3 = 0$ we get $\ln(x) = -\frac{3}{2}$ so $x = e^{-\frac{3}{2}}$.  As with our first derivative analysis, we elect to choose some `log-friendly' test values and note $e^{-2} < e^{-\frac{3}{2}} < e^{0} = 1$.  

We find $f''\left(e^{-2}\right) = 2 \ln \left(e^{-2}\right) + 3 = 2(-2) + 3 = -1 < 0$ and $f''(1) = 2 \ln(1) + 3 = 3 > 0$.
  
\begin{image}
\begin{tikzpicture}[>=Latex, every node/.style={font=\normalsize}]

\def\crit{1}
\def\testL{0}
\def\testR{2}
\def\xmin{-1}
\def\xmax{3.2}

\draw[->] (\xmin,0) -- (\xmax,0);
\draw (-1,0.08) -- (-1,-0.08);
\draw (\crit,0.08) -- (\crit,-0.08);
\node[below] at (\crit,-0.12) {$e^{-3/2}$};
\node[below] at (-1,-0.12) {$0$};
\node[above] at (\crit,0.15) {$0$};

\node[above] at ({(\xmin+\crit)/2},0.15) {$(-)$};
\node[above] at ({(\crit+\xmax-0.6)/2},0.15) {$(+)$};
\node[right] at (\xmax+0.15,0.55) {$f''(x)$};

\draw[->, gray] (\testL,-0.9) -- (\testL,-0.15);
\draw[->, gray] (\testR,-0.9) -- (\testR,-0.15);
\node[right] at (\xmax+0.15,-0.15) {$x$};
\node[below] at (\testL,-0.9) {$e^{-2}$};
\node[below] at (\testR,-0.9) {$1$};

\end{tikzpicture}

\end{image}


Interpreting our sign diagram for $f''(x)$ using Theorem \ref{secondderivatveandgraphs} gives:


\begin{image}
\begin{tikzpicture}[>=Latex, every node/.style={font=\normalsize}]

\def\crit{1}
\def\testL{0}
\def\testR{2}
\def\xmin{-1}
\def\xmax{3.2}

\draw[->] (\xmin,0) -- (\xmax,0);
\draw (-1,0.08) -- (-1,-0.08);
\draw (\crit,0.08) -- (\crit,-0.08);
\node[below] at (\crit,-0.12) {$e^{-3/2}$};
\node[below] at (-1,-0.12) {$0$};
%\node[above] at (\crit,0.15) {$\rightarrow$};

\node[above] at ({(\xmin+\crit)/2},0.15) {C. Down};
\node[above] at ({(\crit+\xmax-0.6)/2},0.12) {C. Up};
\node[right] at (\xmax+0.15,0.55) {$f(x)$};

% \draw[->, gray] (\testL,-0.9) -- (\testL,-0.15);
% \draw[->, gray] (\testR,-0.9) -- (\testR,-0.15);
\node[right] at (\xmax+0.15,-0.15) {$x$};
% \node[below] at (\testL,-0.9) {$e^{-1}$};
% \node[below] at (\testR,-0.9) {$1$};

\end{tikzpicture}

\end{image}

Hence, the graph of $f$ is concave down on $\left(0, e^{-\frac{3}{2}} \right)$ and concave up on $\left(e^{-\frac{3}{2}} , \infty \right)$.


\item  Since the concavity changes at $x = e^{-\frac{3}{2}}$, we have an inflection point at $\left( e^{-\frac{3}{2}}, f \left(e^{-\frac{3}{2}}\right) \right)$.  We find $f \left(e^{-\frac{3}{2}}\right)  = \left(f \left(e^{-\frac{3}{2}}\right) \right)^2 \ln \left(f \left(e^{-\frac{3}{2}}\right) \right) = e^{-3} \left(-\frac{3}{2} \right) = -\frac{3e^{-3}}{2}$.  Our inflection point is $\left( e^{-\frac{3}{2}},  -\frac{3e^{-3}}{2} \right)$.

\item  Using desmos, we confirm our calculations.\footnote{It may require some window adjustment to capture the inflection point.}

% https://www.desmos.com/calculator/yjak0nqwkb
\begin{center}
\desmos{yjak0nqwkb}{800}{600}
\end{center}



\end{enumerate}
\end{explanation}
\end{example}

When determining $\lim_{x \rightarrow 0^{+}} f(x)  = \lim_{x \rightarrow 0^{+}} x^2 \ln(x)$ in Example \ref{logcurvesketchex}, we encountered the indeterminate form `$0 \cdot (-\infty)$.'  This is a similar scenario to what we encountered in the remarks following Example \ref{exponentialcurvesketchingex} in Section \ref{ExponentialEquationsandInequalities}.  In this case, the factor $x^2 \rightarrow 0$ and the factor $\ln(x) \rightarrow -\infty$ and so we have a tug-of-war to see which factor's behavior will win out over the other.  Here, $x^2 \rightarrow 0$ dominates as the table suggests $\lim_{x \rightarrow 0^{+}} x^2 \ln(x) = 0$.  This is indeed the case and, in general, polynomials (and, in general, all positive powers of $x$) dominate logarithms as we'll explore in the Exercise \ref{numericalinvestigationlimitlnxtimesx}.  


%% SKIPPED %% \input{ExercisesforLogarithmicEquationsandInequalities}

\closegraphsfile

\end{document}
